% Introduction (about 7 pages). Written in Stage 7, after the chapters it introduces.
\chapter{Introduction}\label{ch01}

\section{Background and Motivation}\label{sec:ch01-background-and-motivation}

Much of what an organisation knows is split between two kinds of data. Some of it is structured: tables of
products and their parts, vehicles and their recalls, accidents and their findings, each row identified by a key and
joined to others. The rest is text: customer reviews, owners' complaints, accident narratives, reports and minutes.
The two describe the same things --- a review names the product and the part that failed; a complaint names the
vehicle a recall covers --- but they are rarely connected. A question that needs both, such as \emph{which suppliers
provide the parts that reviewers report as defective}, requires someone to read the text, recognise what it refers
to, find the matching records, and follow the records' relations.

Knowledge graphs offer a way to connect the two: entities become nodes, relations become edges, and both tables and
text can be mapped into one structure that can be navigated and queried~\cite{hogan2021knowledge}. Building such a
graph has traditionally required a hand-designed schema, extraction rules written for one domain, and careful entity
resolution. Large language models (LLMs) change this: they can propose a schema for data they have never seen, extract
facts from text, and judge whether two names refer to the same thing~\cite{pan2024unifying,zhu2024llms}. They can
also answer questions over retrieved evidence~\cite{lewis2020retrieval} and act as agents that call
tools~\cite{yao2023react}. But they are fallible in ways that matter for a knowledge graph: an extracted fact may not
be in the text, a proposed schema may not cover what the text says, and two different things may be merged because
their names are similar.

This thesis builds and evaluates a system that uses LLMs for what they do well --- proposing --- and code for what it
does well --- checking and computing exactly. It turns mixed tables and text into a knowledge graph and answers
natural-language questions over it, and it measures, at each step, whether the result is correct.

\section{Problem Statement}\label{sec:ch01-problem-statement}

Two problems are addressed.

\paragraph{Construction.} Given a folder of tables and documents and a goal stated in one sentence, build
automatically, without code or settings written for the dataset, a graph in which the things the data talks about can
be found and from which the evidence about each of them can be reached, with every element traceable to the record or
passage that supports it. The hard part is not extraction alone but \emph{identity}: deciding which mention in which
document refers to which record, which mentions denote the same individual, and which must be kept apart.

\paragraph{Querying.} Given such a graph, its text and its tables, answer questions in natural language --- including
questions that count, filter or join records, follow several hops, or collect every mention in the text --- with
citations to the evidence, and recognise questions that should not be answered: unclear, out of scope, without
evidence, or asking to change the data.

Behind both lies a design question that the thesis treats as a research question rather than an assumption:
\emph{what should the graph contain?} A graph can store every fact extracted from the text, or it can serve as an index
that links the text to the records and things it mentions and leaves the text to be read when a question is asked.

\section{Objectives}\label{sec:ch01-objectives}

\begin{enumerate}[nosep]
\item Build a domain-neutral pipeline that constructs a provenance-aware, entity-centric graph from tables and text,
with every model output validated by code.
\item Compare, by measurement, graph representations that store extracted facts with one that anchors identity and
evidence.
\item Build a query engine that answers natural-language questions over the graph, the text and the tables, with
citations checked by code and explicit handling of questions that should not be answered.
\item Evaluate both on several datasets with blind reference sets, two kinds of score and paired statistical tests,
and report every number with its provenance.
\end{enumerate}

\section{Research Questions}\label{sec:ch01-research-questions}

\begin{description}[style=nextline]
\item[RQ1 (construction)] To what extent can an LLM-proposed, code-verified pipeline turn heterogeneous tables and
text into an entity-centric graph whose identity, record links and provenance are correct, without domain-specific
code?
\item[RQ2 (representation)] For locating evidence and supporting question answering, what are the measured trade-offs
between storing extracted facts as the graph's backbone and using the graph as an entity-centric index over records and
text?
\item[RQ3 (query architecture)] How accurately can a tool-using agent that combines graph navigation, text retrieval
and typed structured queries answer natural-language questions over such data, with verifiable citations, compared
with retrieval-only baselines?
\item[RQ4 (design factors)] How do retrieval choices, tool expressiveness and the model's reasoning budget affect
accuracy, cost and latency?
\end{description}

RQ1 and RQ2 are answered in \cref{ch05}, RQ3 and RQ4 in \cref{ch07}; \cref{ch09} brings the answers together.

\section{Scope and Non-Goals}\label{sec:ch01-scope-and-non-goals}

\begin{itemize}[nosep]
\item \emph{Mixed data only.} Every dataset combines tables and text; text-only and table-only settings are not
studied.
\item \emph{Small corpora.} The four datasets have tens of documents and at most a few hundred rows each, so that every
judged number could be checked against the sources in full. Scale is not studied.
\item \emph{Short documents in English.} Long narratives, scientific papers, images, figures and other languages are
out of scope.
\item \emph{Questions about the data.} The engine answers questions the data can answer; it does not summarise a whole
corpus, and it never changes the data.
\item \emph{One model family per role.} The system uses Gemini models; the reference answers and judgements come from
Claude. Comparing model families is not a goal.
\end{itemize}

\section{Contributions}\label{sec:ch01-contributions}

The thesis makes six contributions; \cref{sec:ch02-positioning-of-this-work} positions them against related work.
\begin{enumerate}[nosep]
\item A domain-neutral construction pipeline for mixed tables and text in which a language model proposes and code
decides, including identity resolution that asks a model only about near misses and accepts its answer only with
verified quotations.
\item The \emph{anchor graph}: a representation that uses the graph as an index of identity and evidence, with identity
stored as reversible, evidenced edges, a rule that a shared node finds passages but never proves a fact, and a
navigation contract.
\item An evaluation of a graph \emph{as an index}, by ten criteria with hard gates fixed before measurement.
\item A measured design history comparing four graph representations on the same data.
\item A tool-using query engine with typed structured queries, code-checked citations and answer outcomes, a
deterministic screen and read-only access, evaluated on blind question sets.
\item A reproducible, low-cost evaluation practice in which every reported number is traced to its run, its record and
its commit.
\end{enumerate}
The techniques the system combines --- BM25 and dense retrieval, Reciprocal Rank Fusion, reranking, the ReAct pattern,
blocking for entity resolution, McNemar's test --- are established; the contribution lies in how they are combined and
constrained for mixed tables and text, in the evidence on which representation serves this purpose, and in the
evaluation itself.

\section{Thesis Organisation}\label{sec:ch01-thesis-organisation}

\Cref{ch02} reviews the background: knowledge graphs, extraction, entity resolution, provenance, retrieval,
retrieval-augmented generation, agents and evaluation. \Cref{ch03} states the research method, the datasets, the
reference sets, the scores, the metrics and the statistics. \Cref{ch04} describes the builder and the evolution from a
graph of facts to the anchor graph, and \cref{ch05} evaluates it. \Cref{ch06} describes the query engine and
\cref{ch07} evaluates it. \Cref{ch08} looks at the integrated system, the division of work between its parts, and the
threats to validity. \Cref{ch09} answers the research questions and outlines future work. The appendices give the
prompts, protocols, complete results, the excluded dataset's measurements, and a guide to reproducing the results.
